\chapter{Standard Template Library Core}
\label{ch:stl-core}

\textit{Generic containers, iterators, and algorithms --- the vocabulary every C++ program uses every day.}

The \textbf{Standard Template Library (STL)} is the most important library in C++. It delivers production-quality implementations of the most common data structures (containers), traversal abstractions (iterators), and algorithms --- all expressed as templates, so they work with any element type. This chapter covers the C++98/03 STL in full: the nine containers, the five iterator categories, the algorithm header, \texttt{std::string}, streams, and functors.

%% ============================================================
\section{STL Architecture}
\label{sec:stl-architecture}

The STL is organised around four cooperating components:

\begin{lstlisting}[language={},caption={STL component tree}]
          STL (Standard Template Library)
                       |
   -----------------------------------------
   |              |            |           |
CONTAINERS     ITERATORS    ALGORITHMS   FUNCTORS
   |              |            |           |
Sequence       Input        Searching    Predicates
Associative    Output       Sorting      Comparators
Adapters       Forward      Modifying
               Bidir        Numeric
               Random
\end{lstlisting}

\textbf{Containers} store data. \textbf{Iterators} point into containers and define ranges. \textbf{Algorithms} operate on iterator ranges, blind to the underlying container. \textbf{Functors} parametrise algorithm behaviour without calling convention overhead.

Key advantages:
\begin{itemize}
    \item \textbf{Generic programming}: algorithms written once work with any container of any element type.
    \item \textbf{Separation of concerns}: swapping from \texttt{vector} to \texttt{list} changes only the container type, not the algorithm.
    \item \textbf{Performance}: implementations are heavily optimised; the abstractions typically compile to zero overhead.
\end{itemize}

%% ============================================================
\section{Container Characteristics at a Glance}
\label{sec:container-overview}

\begin{center}
\begin{tabular}{|l|l|l|l|l|l|}
\hline
\textbf{Container} & \textbf{Type} & \textbf{Insert} & \textbf{Search} & \textbf{Random Access} & \textbf{Memory} \\
\hline
\texttt{vector}   & Sequence    & O(1) back†   & O(n)      & O(1)  & Contiguous \\
\texttt{deque}    & Sequence    & O(1) ends    & O(n)      & O(1)  & Block chunks \\
\texttt{list}     & Sequence    & O(1)‡        & O(n)      & None  & Scattered \\
\texttt{map/set}  & Associative & O(log n)     & O(log n)  & None  & Red-Black Tree \\
\texttt{stack}    & Adapter     & O(1)         & ---       & None  & Underlying \\
\texttt{queue}    & Adapter     & O(1)         & ---       & None  & Underlying \\
\texttt{priority\_queue} & Adapter & O(log n)  & ---       & None  & Heap (vector) \\
\hline
\end{tabular}
\end{center}

†Amortised; reallocation is O(n). ‡Given an iterator to the target node.

%% ============================================================
\section{\texttt{std::vector} --- The Dynamic Array}
\label{sec:vector}

\textbf{\texttt{std::vector}} is the default container. It stores elements contiguously in heap memory, growing automatically. Use it unless you have a specific reason not to.

\begin{lstlisting}[language=C++, caption={Declaring and initialising vectors}]
#include <vector>
using namespace std;

vector<int> v1;               // Empty
vector<int> v2(10);           // 10 elements, value-initialised to 0
vector<int> v3(5, 42);        // 5 elements, all 42
vector<int> v4(v3);           // Copy

int arr[] = {1, 2, 3, 4, 5};
vector<int> v5(arr, arr + 5); // Construct from C array
\end{lstlisting}

\subsection{Element Access}

\begin{lstlisting}[language=C++, caption={Vector element access}]
vector<int> v;
v.push_back(10);
v.push_back(20);

v[0];       // No bounds check
v.at(1);    // Throws std::out_of_range if out of range
v.front();  // 10
v.back();   // 20
\end{lstlisting}

\subsection{Modifying a Vector}

\begin{lstlisting}[language=C++, caption={Common vector mutations}]
vector<int> v;
v.push_back(10); v.push_back(20); v.push_back(30); // {10, 20, 30}
v.insert(v.begin() + 1, 15);  // {10, 15, 20, 30}
v.pop_back();                  // {10, 15, 20}
v.erase(v.begin() + 1);        // {10, 20}
v.clear();                     // {}
\end{lstlisting}

\subsection{Size and Capacity}

\begin{lstlisting}[language=C++, caption={size vs capacity}]
vector<int> v(10);
v.size();      // 10 -- elements actually stored
v.capacity();  // >= 10 -- allocated slots
v.reserve(100); // Allocate space for 100, no new elements
v.resize(20);   // Add default-initialised elements to size 20
v.resize(5);    // Truncate to 5 elements
\end{lstlisting}

Always \texttt{reserve()} if you know the final size --- reallocation copies every element and invalidates all iterators.

\subsection{Iterating}

\begin{lstlisting}[language=C++, caption={Three iteration patterns}]
#include <vector>
#include <iostream>
using namespace std;

int main() {
    vector<int> v;
    v.push_back(10); v.push_back(20); v.push_back(30);

    for (size_t i = 0; i < v.size(); ++i)           // Index
        cout << v[i] << " ";

    for (vector<int>::iterator it = v.begin(); it != v.end(); ++it) // Iterator
        cout << *it << " ";

    for (vector<int>::reverse_iterator rit = v.rbegin(); rit != v.rend(); ++rit)
        cout << *rit << " ";
    return 0;
}
\end{lstlisting}

%% ============================================================
\section{\texttt{std::deque} --- Double-Ended Queue}
\label{sec:deque}

\begin{lstlisting}[language=C++, caption={deque basics}]
#include <deque>
#include <iostream>
using namespace std;

int main() {
    deque<int> dq;
    dq.push_back(10);
    dq.push_front(5);   // {5, 10}
    dq.pop_back();      // {5}
    dq.pop_front();     // {}
    dq.push_back(100);
    cout << dq[0] << "\n"; // Random access
    return 0;
}
\end{lstlisting}

%% ============================================================
\section{\texttt{std::list} --- Doubly Linked List}
\label{sec:list}

\begin{lstlisting}[language=C++, caption={list operations}]
#include <list>
using namespace std;

int main() {
    list<int> lst;
    lst.push_back(10); lst.push_front(5); // {5, 10}

    list<int>::iterator it = lst.begin();
    ++it;
    lst.insert(it, 7); // {5, 7, 10}

    lst.remove(7);     // {5, 10}
    lst.push_back(10);
    lst.unique();      // Remove consecutive duplicates -> {5, 10}
    lst.sort();        // std::sort won't work; use member sort()
    return 0;
}
\end{lstlisting}

%% ============================================================
\section{\texttt{std::map} --- Sorted Key-Value Store}
\label{sec:map}

\begin{lstlisting}[language=C++, caption={map operations}]
#include <map>
#include <string>
#include <iostream>
using namespace std;

int main() {
    map<string, int> ages;
    ages["Alice"] = 30;
    ages["Bob"]   = 25;
    ages.insert(make_pair("Charlie", 28));

    cout << ages["Alice"] << "\n"; // WARNING: inserts default if absent!

    map<string, int>::iterator it = ages.find("Bob");
    if (it != ages.end())
        cout << "Bob is " << it->second << " years old.\n";

    for (it = ages.begin(); it != ages.end(); ++it)
        cout << it->first << ": " << it->second << "\n";
    return 0;
}
\end{lstlisting}

\begin{lstlisting}[language=C++, caption={map with custom comparator}]
#include <map>
#include <string>
#include <cstring>

struct CaseInsensitive {
    bool operator()(const std::string& a, const std::string& b) const {
        return strcasecmp(a.c_str(), b.c_str()) < 0;
    }
};
std::map<std::string, int, CaseInsensitive> ci_map;
\end{lstlisting}

%% ============================================================
\section{\texttt{std::set}, \texttt{std::multimap}, \texttt{std::multiset}}
\label{sec:set-multi}

\begin{lstlisting}[language=C++, caption={set and multimap}]
#include <set>
#include <map>
#include <string>
#include <iostream>
using namespace std;

int main() {
    // set -- sorted, unique
    set<int> s;
    s.insert(10); s.insert(5); s.insert(10); // {5, 10}
    if (s.find(5) != s.end()) cout << "5 found\n";
    s.erase(s.begin());

    // multimap -- allows duplicate keys
    multimap<string, int> scores;
    scores.insert(make_pair("Alice", 90));
    scores.insert(make_pair("Alice", 85));

    typedef multimap<string, int>::iterator MMit;
    pair<MMit, MMit> range = scores.equal_range("Alice");
    for (MMit it = range.first; it != range.second; ++it)
        cout << it->second << " ";
    return 0;
}
\end{lstlisting}

%% ============================================================
\section{Container Adapters: \texttt{stack}, \texttt{queue}, \texttt{priority\_queue}}
\label{sec:adapters}

\begin{lstlisting}[language=C++, caption={Container adapters}]
#include <stack>
#include <queue>
#include <vector>
#include <functional>
#include <iostream>
using namespace std;

int main() {
    // LIFO stack
    stack<int> st;
    st.push(10); st.push(20);
    cout << st.top() << "\n"; // 20
    st.pop();

    // FIFO queue
    queue<int> q;
    q.push(10); q.push(20);
    cout << q.front() << "\n"; // 10
    q.pop();

    // Max-heap
    priority_queue<int> pq;
    pq.push(10); pq.push(30); pq.push(20);
    cout << pq.top() << "\n"; // 30

    // Min-heap
    priority_queue<int, vector<int>, greater<int> > min_pq;
    min_pq.push(10); min_pq.push(30);
    cout << min_pq.top() << "\n"; // 10

    return 0;
}
\end{lstlisting}

%% ============================================================
\section{Iterators}
\label{sec:iterators}

\begin{center}
\begin{tabular}{|l|l|l|}
\hline
\textbf{Category} & \textbf{Example containers} & \textbf{Extra operations} \\
\hline
Input            & \texttt{istream\_iterator} & \texttt{*it}, \texttt{++it} (single pass, read) \\
Output           & \texttt{ostream\_iterator} & \texttt{*it=v}, \texttt{++it} (single pass, write) \\
Forward          & \texttt{slist} (non-std)   & Input + multi-pass \\
Bidirectional    & \texttt{list}, \texttt{map}, \texttt{set} & Forward + \texttt{-{}-it} \\
Random Access    & \texttt{vector}, \texttt{deque} & Bidir + \texttt{it+n}, \texttt{it-it2}, \texttt{it[n]} \\
\hline
\end{tabular}
\end{center}

\begin{lstlisting}[language=C++, caption={advance and distance}]
#include <vector>
#include <iterator>
#include <iostream>
using namespace std;

int main() {
    vector<int> v;
    v.push_back(1); v.push_back(2); v.push_back(3); v.push_back(4); v.push_back(5);

    vector<int>::iterator it = v.begin();
    advance(it, 2);                            // Now points to element 3
    cout << distance(v.begin(), it) << "\n";  // 2
    cout << *it << "\n";                       // 3
    ++it; // Prefer ++it over it++ for all iterators
    return 0;
}
\end{lstlisting}

%% ============================================================
\section{Algorithms}
\label{sec:algorithms}

Algorithms operate on half-open ranges \texttt{[first, last)} --- \texttt{last} is past-the-end and must never be dereferenced.

\begin{lstlisting}[language=C++, caption={Non-modifying, modifying, and sorting algorithms}]
#include <algorithm>
#include <vector>
#include <iostream>
using namespace std;

bool descending(int a, int b) { return a > b; }

int main() {
    int arr[] = {5, 2, 9, 1, 5, 6};
    vector<int> v(arr, arr + 6);

    // find / count
    vector<int>::iterator it = find(v.begin(), v.end(), 9);
    cout << count(v.begin(), v.end(), 5) << "\n"; // 2

    // sort ascending, then descending
    sort(v.begin(), v.end());
    sort(v.begin(), v.end(), descending);

    // binary search (range must be sorted)
    sort(v.begin(), v.end());
    cout << binary_search(v.begin(), v.end(), 5) << "\n"; // 1

    vector<int>::iterator lo = lower_bound(v.begin(), v.end(), 5);
    vector<int>::iterator hi = upper_bound(v.begin(), v.end(), 5);
    cout << "Count of 5: " << (hi - lo) << "\n";

    // reverse, fill, copy
    reverse(v.begin(), v.end());
    fill(v.begin(), v.end(), 0);
    return 0;
}
\end{lstlisting}

\begin{lstlisting}[language=C++, caption={Numeric algorithms}]
#include <numeric>
#include <vector>
#include <iostream>
using namespace std;

int main() {
    int arr[] = {1, 2, 3, 4, 5};
    vector<int> v(arr, arr + 5);

    int total = accumulate(v.begin(), v.end(), 0); // 15
    int dot   = inner_product(v.begin(), v.end(), v.begin(), 0); // 55

    cout << total << "\n" << dot << "\n";
    return 0;
}
\end{lstlisting}

%% ============================================================
\section{\texttt{std::string}}
\label{sec:string}

\begin{lstlisting}[language=C++, caption={std::string fundamentals}]
#include <string>
#include <iostream>
using namespace std;

int main() {
    string s = "Hello";
    s += " World";               // "Hello World"

    string sub = s.substr(0, 5); // "Hello"

    size_t pos = s.find("World");
    if (pos != string::npos)
        cout << "Found at " << pos << "\n"; // 6

    s.replace(6, 5, "C++");    // "Hello C++"
    s.insert(5, " there");     // "Hello there C++"
    s.erase(5, 6);             // "Hello C++"

    const char* cstr = s.c_str(); // null-terminated const pointer
    return 0;
}
\end{lstlisting}

%% ============================================================
\section{Streams and String Streams}
\label{sec:streams}

\begin{lstlisting}[language=C++, caption={Basic file I/O}]
#include <fstream>
#include <string>
#include <iostream>
using namespace std;

int main() {
    ofstream out("test.txt");
    if (out) { out << "Hello file\n" << 42 << "\n"; }
    out.close();

    ifstream in("test.txt");
    string line;
    while (getline(in, line)) cout << line << "\n";
    in.close();
    return 0;
}
\end{lstlisting}

\begin{lstlisting}[language=C++, caption={stringstream for type conversion (C++98/03 idiom)}]
#include <sstream>
#include <string>
#include <iostream>
using namespace std;

int main() {
    int x = 42;
    stringstream ss;
    ss << x;
    string s = ss.str(); // "42"

    string s2 = "100";
    stringstream ss2(s2);
    int y;
    ss2 >> y; // y == 100

    cout << s << " " << y << "\n";
    return 0;
}
\end{lstlisting}

%% ============================================================
\section{Functors (Function Objects)}
\label{sec:functors}

\begin{lstlisting}[language=C++, caption={Custom functor with state}]
#include <vector>
#include <algorithm>
#include <iostream>
using namespace std;

struct AddX {
    int x;
    explicit AddX(int val) : x(val) {}
    int operator()(int y) const { return x + y; }
};

int main() {
    int arr[] = {1, 2, 3, 4, 5};
    vector<int> v(arr, arr + 5);

    transform(v.begin(), v.end(), v.begin(), AddX(10));
    // v is now {11, 12, 13, 14, 15}

    for (vector<int>::iterator it = v.begin(); it != v.end(); ++it)
        cout << *it << " ";
    cout << "\n";
    return 0;
}
\end{lstlisting}

\begin{lstlisting}[language=C++, caption={Standard functors from <functional>}]
#include <algorithm>
#include <functional>
#include <vector>
using namespace std;

int main() {
    int arr[] = {5, 1, 3, 9, 2};
    vector<int> v(arr, arr + 5);
    sort(v.begin(), v.end(), greater<int>()); // {9, 5, 3, 2, 1}
    return 0;
}
\end{lstlisting}

%% ============================================================
\section{Advanced String Operations}
\label{sec:string-advanced}

\begin{lstlisting}[language=C++, caption={String manipulation}]
#include <iostream>
#include <string>
#include <algorithm>
using namespace std;

int main() {
    string s = "Hello World";
    size_t pos = s.find("World");
    s.replace(6, 5, "C++");   // "Hello C++"
    s.insert(5, " there");    // "Hello there C++"
    s.erase(5, 6);            // "Hello C++"
    cout << s.substr(0, 5) << "\n"; // "Hello"
    reverse(s.begin(), s.end());
    cout << s << "\n"; // "++C olleH"
    return 0;
}
\end{lstlisting}

\begin{lstlisting}[language=C++, caption={C++98 string tokenisation}]
#include <iostream>
#include <string>
#include <sstream>
using namespace std;

int main() {
    string line = "apple,banana,orange,grape";
    stringstream ss(line);
    string token;
    while (getline(ss, token, ','))
        cout << token << "\n";
    return 0;
}
\end{lstlisting}

%% ============================================================
\section{Professional Insights: STL Internals and Pitfalls}
\label{sec:stl-pro-insights}

\subsection{Container Selection}

\begin{center}
\begin{tabular}{|l|l|}
\hline
\textbf{Need} & \textbf{Best container} \\
\hline
Fast random access, cache-friendly iteration & \texttt{vector} \\
Fast push/pop at both ends                   & \texttt{deque} \\
Fast insert/erase with stable iterators      & \texttt{list} \\
Key-based lookup, sorted traversal           & \texttt{map} / \texttt{set} \\
Priority queue                               & \texttt{priority\_queue} \\
Bounded LIFO/FIFO                            & \texttt{stack} / \texttt{queue} \\
\hline
\end{tabular}
\end{center}

\subsection{Iterator Invalidation}

Misusing invalidated iterators is undefined behaviour. Key rules:

\begin{itemize}
    \item \textbf{vector}: Any insertion that causes reallocation invalidates all iterators. Insertions at or after point of change also invalidate.
    \item \textbf{deque}: Any insertion or erasure invalidates all iterators.
    \item \textbf{list / map / set}: Insertion never invalidates existing iterators. Erasure only invalidates the erased node's iterator.
\end{itemize}

\subsection{\texttt{vector<bool>} Is Not a Regular Vector}

\texttt{std::vector<bool>} packs bits one per bit. \texttt{operator[]} returns a \textbf{proxy object}, not a \texttt{bool\&}. This breaks generic code. Prefer \texttt{vector<char>} or \texttt{std::bitset<N>}.

\subsection{Prefer \texttt{++it} Over \texttt{it++}}

Postfix \texttt{it++} creates a temporary iterator copy. Prefix \texttt{++it} advances in place. Always use prefix in loops.

\subsection{Strict Weak Ordering for Comparators}

Any comparator passed to \texttt{sort}, \texttt{map}, or \texttt{priority\_queue} must be irreflexive, asymmetric, and transitive. Using \texttt{<=} instead of \texttt{<} is undefined behaviour.

\subsection{The \texttt{equal\_range} Pattern}

\begin{lstlisting}[language=C++, caption={equal\_range for multi-element lookup}]
#include <algorithm>
#include <vector>
#include <iostream>
using namespace std;

int main() {
    int arr[] = {1, 2, 5, 5, 5, 9};
    vector<int> v(arr, arr + 6);

    pair<vector<int>::iterator, vector<int>::iterator> range
        = equal_range(v.begin(), v.end(), 5);

    cout << "Count of 5: " << (range.second - range.first) << "\n"; // 3
    return 0;
}
\end{lstlisting}
